\documentclass[../../main.tex]{subfiles}
\begin{document}
\begin{oframed}
  \subfile{problem}
\end{oframed}

{\bf [解答]}

\begin{figure}[htb]
  \centering
  \begin{tikzpicture}[x=1.25cm, y=0.72cm, >=stealth, every node/.style={font=\small}]
    % 左パネル: y 軸に平行な直線上の 3 点
    \begin{scope}
      \node[font=\footnotesize] at (1.9, 6.6) {(a) $y$軸に平行な直線（$3\,{}_{n}C_{3}$ 組）};
      \draw[->, thick, gray!80!black] (0, 0) -- (3.8, 0) node[right, black] {$x$};
      \draw[->, thick, gray!80!black] (0, 0) -- (0, 5.7) node[above, black] {$y$};
      \foreach \xi in {1, 2, 3} {
        \node[below=2pt, font=\scriptsize] at (\xi, 0) {$\xi$};
      }
      \foreach \yi in {1, 2, 3, 4, 5} {
        \node[left=2pt, font=\scriptsize] at (0, \yi) {$\yi$};
      }
      \draw[line width=1.2pt, blue!85!black] (2, 0.5) -- (2, 5.4);
      \node[above, blue!85!black, font=\footnotesize] at (2, 5.4) {$x=2$};
      \foreach \xi in {1, 2, 3} {
        \foreach \yi in {1, 2, 3, 4, 5} {
          \node[circle, fill=gray!70!black, inner sep=1.2pt] at (\xi, \yi) {};
        }
      }
      \foreach \yi in {1, 3, 4} {
        \node[circle, fill=red!85!black, inner sep=2.0pt] at (2, \yi) {};
      }
    \end{scope}

    % 右パネル: それ以外の直線上の 3 点
    \begin{scope}[xshift=6.8cm]
      \node[font=\footnotesize] at (2.0, 6.6) {(b) それ以外の直線（$2y_2 = y_1 + y_3$）};
      \draw[->, thick, gray!80!black] (0, 0) -- (3.9, 0) node[right, black] {$x$};
      \draw[->, thick, gray!80!black] (0, 0) -- (0, 5.7) node[above, black] {$y$};
      \foreach \xi in {1, 2, 3} {
        \node[below=2pt, font=\scriptsize] at (\xi, 0) {$\xi$};
      }
      \foreach \yi in {1, 2, 3, 4, 5} {
        \node[left=2pt, font=\scriptsize] at (0, \yi) {$\yi$};
      }
      \draw[densely dotted, thick, gray!70!black] (0, 1) -- (1, 1);
      \draw[densely dotted, thick, gray!70!black] (0, 3) -- (2, 3);
      \draw[densely dotted, thick, gray!70!black] (0, 5) -- (3, 5);
      \draw[line width=1.2pt, blue!85!black] (0.7, 0.4) -- (3.3, 5.6);
      \foreach \xi in {1, 2, 3} {
        \foreach \yi in {1, 2, 3, 4, 5} {
          \node[circle, fill=gray!70!black, inner sep=1.2pt] at (\xi, \yi) {};
        }
      }
      \node[circle, fill=red!85!black, inner sep=2.0pt] at (1, 1) {};
      \node[circle, fill=red!85!black, inner sep=2.0pt] at (2, 3) {};
      \node[circle, fill=red!85!black, inner sep=2.0pt] at (3, 5) {};
      \node[below right=2pt, red!85!black, font=\scriptsize, fill=white, inner sep=0.8pt] at (1, 1) {$(1, y_1)$};
      \node[below right=2pt, red!85!black, font=\scriptsize, fill=white, inner sep=0.8pt] at (2, 3) {$(2, y_2)$};
      \node[below right=2pt, red!85!black, font=\scriptsize, fill=white, inner sep=0.8pt] at (3, 5) {$(3, y_3)$};
    \end{scope}
  \end{tikzpicture}
  \caption{選んだ3点が一直線上に並ぶ配置（例：$n=5$）}
\end{figure}

与えられた集合
\begin{align}
  \{ (x, \, y) \, | \, x, \, y \text{は}
  1\leqq x \leqq3, \, 1\leqq y \leqq n \text{を満たす整数} \}
\end{align}
を考える．まず$3$点を選ぶ場合の数は
\begin{align}
  {}_{3n}C_{3} \quad \text{通り} \label{eq:1}
\end{align}
で同様に確からしい．

選んだ相異なる3点が三角形を作らないのは$3$点が一直線上にあるときである．
まず直線が$y$軸平行な場合は$x=1,x=2,x=3$の$3$本があり，それぞれ点の選択肢は$n$個だから全部で
\begin{align}
  3{}_{n}C_{3} \quad \text{組} \label{eq:2}
\end{align}
ある．

それ以外の直線の時，$x$座標$1,2,3$からそれぞれ$1$点ずつが選ばれるから，これらを$(1,y_1),(2,y_2),(3,y_3)$とおく．するとこれらが一直線上に並ぶ条件は
\begin{align}
  y_3 - y_2 = y_2 - y_1
  \quad\therefore
  2y_2 = y_1 + y_3 \label{eq:3}
\end{align}
である．左辺が偶数だから右辺も偶数である必要があり，よって$y_1,y_3$の偶奇は等しい．
逆に等しい偶奇で$y_1,y_3$を選べば\cref{eq:3}に応じて$1 \le y_2 \le n$が唯一定まる．

\paragraph{(1)}

$n=5$のとき，同じ偶奇の$(y_1,y_3)$は$3^2+2^2=13$組である．
よって求める確率は排反事象を考えて\cref{eq:1,eq:2}より
\begin{align}
  p_5
  = 1-\dfrac{3{}_{5}C_{3}+13}{{}_{15}C_{3}}
  =\dfrac{412}{455}
\end{align}
である．

\paragraph{(2)}

$n=2m$のとき同じ偶奇の$(y_1,y_3)$は$2m^2$組である．したがって求める確率は排反事象を考えて\cref{eq:1,eq:2}より
\begin{align}
  p_{2m}
  &=1-\dfrac{3{}_{2m}C_{3}+2m^2}{{}_{6m}C_{3}}\\
  &=1-\dfrac{2(2m^2-2m+1)}{(6m-1)(6m-2)} \\
  &=\dfrac{m(16m-7)}{(6m-1)(3m-1)}
\end{align}
である．

\end{document}
